\hypertarget{chapter-19-c14-standard-library-enhancements}{%
\section{CHAPTER 19: C14 STANDARD LIBRARY
ENHANCEMENTS}\label{chapter-19-c14-standard-library-enhancements}}

\hypertarget{c14-standard-library-enhancements}{%
\section{C++14 STANDARD LIBRARY
ENHANCEMENTS}\label{c14-standard-library-enhancements}}

The C++14 library updates were targeted at consistency and fixing
omissions from C++11. Most notably, it finally gave us
\passthrough{\lstinline!std::make\_unique!} and introduced the first
standardized reader-writer lock.

\hypertarget{stdmake_unique-completing-the-set}{%
\subsection{\texorpdfstring{1. \texttt{std::make\_unique}: Completing
the
Set}{1. std::make\_unique: Completing the Set}}\label{stdmake_unique-completing-the-set}}

In C++11, we had \passthrough{\lstinline!std::make\_shared!} but no
\passthrough{\lstinline!std::make\_unique!}. This was a strange omission
that forced developers to use \passthrough{\lstinline!new!} for unique
pointers.

\hypertarget{why-use-make_unique}{%
\subsubsection{\texorpdfstring{1.1 Why use
\texttt{make\_unique}?}{1.1 Why use make\_unique?}}\label{why-use-make_unique}}

\begin{enumerate}
\def\labelenumi{\arabic{enumi}.}
\tightlist
\item
  \textbf{Exception Safety:} Prevents memory leaks in complex
  expressions where multiple allocations occur.
\item
  \textbf{No \passthrough{\lstinline!new!} Keyword:} Keeps code clean
  and adheres to the ``No Raw New/Delete'' modern C++ philosophy.
\item
  \textbf{Efficiency:} While it doesn't offer the control-block
  optimization of \passthrough{\lstinline!make\_shared!}, it is the
  standard way to construct unique pointers.
\end{enumerate}

\begin{lstlisting}[language={C++}]
// BAD: Potential leak if foo() throws
process(std::unique_ptr<T>(new T()), foo());

// GOOD: Exception safe
process(std::make_unique<T>(), foo());
\end{lstlisting}

\hypertarget{stdexchange-move-semantics-utility}{%
\subsection{\texorpdfstring{2. \texttt{std::exchange}: Move Semantics
Utility}{2. std::exchange: Move Semantics Utility}}\label{stdexchange-move-semantics-utility}}

\passthrough{\lstinline!std::exchange!} replaces the value of an object
with a new value and returns the old value. It is particularly useful
for implementing move constructors and move assignment operators.

\begin{lstlisting}[language={C++}]
struct Node {
    int* data;
    Node(Node&& other) noexcept
        : data(std::exchange(other.data, nullptr)) {}

    Node& operator=(Node&& other) noexcept {
        if (this != &other) {
            delete data;
            data = std::exchange(other.data, nullptr);
        }
        return *this;
    }
};
\end{lstlisting}

\hypertarget{stdshared_timed_mutex-reader-writer-lock}{%
\subsection{\texorpdfstring{3. \texttt{std::shared\_timed\_mutex}
(Reader-Writer
Lock)}{3. std::shared\_timed\_mutex (Reader-Writer Lock)}}\label{stdshared_timed_mutex-reader-writer-lock}}

One of the most requested features: a mutex that allows multiple readers
OR one writer.

\hypertarget{performance-considerations}{%
\subsubsection{3.1 Performance
Considerations}\label{performance-considerations}}

Use a shared mutex when: - Reads are frequent and cheap. - Writes are
infrequent and expensive.

\begin{lstlisting}[language={C++}]
#include <shared_mutex>
#include <map>

class ThreadSafeMap {
    std::map<int, std::string> data;
    mutable std::shared_timed_mutex mtx;

public:
    std::string get(int key) const {
        std::shared_lock lock(mtx); // Shared lock (Read)
        return data.at(key);
    }

    void set(int key, std::string val) {
        std::unique_lock lock(mtx); // Exclusive lock (Write)
        data[key] = std::move(val);
    }
};
\end{lstlisting}

\hypertarget{stdquoted-stream-parsing-hero}{%
\subsection{\texorpdfstring{4. \texttt{std::quoted}: Stream Parsing
Hero}{4. std::quoted: Stream Parsing Hero}}\label{stdquoted-stream-parsing-hero}}

Parsing CSVs or logs with quoted strings used to be a nightmare of
manual character escaping. \passthrough{\lstinline!std::quoted!} handles
this automatically.

\begin{lstlisting}[language={C++}]
#include <iomanip>
#include <sstream>

void test_quoted() {
    std::stringstream ss;
    std::string s = "Hello \"C++14\" World";

    ss << std::quoted(s);
    // ss now contains: "Hello \"C++14\" World" (with quotes and escaping)

    std::string output;
    ss >> std::quoted(output);
    // output now contains original string: Hello "C++14" World
}
\end{lstlisting}

\hypertarget{volume-03-godhood-summary}{%
\section{VOLUME 03: GODHOOD SUMMARY}\label{volume-03-godhood-summary}}

C++14 was the release where \textbf{Modern C++ became ``Fluid.''}

\begin{enumerate}
\def\labelenumi{\arabic{enumi}.}
\tightlist
\item
  \textbf{Constexpr is King:} Logic migrated from runtime to
  compile-time. If it doesn't involve I/O or dynamic allocation, it
  should probably be \passthrough{\lstinline!constexpr!}.
\item
  \textbf{Lambdas are Complete:} With Generic Lambdas and Init-Capture,
  lambdas are now the preferred tool for almost all local logic,
  closures, and callback patterns.
\item
  \textbf{Standard Consistency:}
  \passthrough{\lstinline!std::make\_unique!} and
  \passthrough{\lstinline!\_t/\_v!} aliases removed the ``boilerplate
  friction'' that made C++11 feel verbose.
\end{enumerate}

\textbf{The Golden Rule of C++14:} If you find yourself writing a manual
loop in a template or a manual move in a constructor, check if a C++14
utility (\passthrough{\lstinline!integer\_sequence!},
\passthrough{\lstinline!std::exchange!}) can do it for you in one line.

\hypertarget{volume-04-simplification-modernization-c17}{%
\section{VOLUME 04 SIMPLIFICATION MODERNIZATION
C17}\label{volume-04-simplification-modernization-c17}}
